%%%%%%%%%%%%%%%%%%%%%%%%%%%%%%%%%
%
%       EXERCÍCIO
%
%%%%%%%%%%%%%%%%%%%%%%%%%%%%%%%%%

\definevalorquestao

\begin{exercicioBanco}[\valorquestao]
A tabela a seguir apresenta a distribuição de frequências relativas aos tempos de atendimento (em minutos) registrados em um posto de saúde:

\begin{center}
\begin{tabular}{|c|c|}
\hline
Tempo (min) & Frequência ($f_i$) \\ \hline
$0 \vdash 10$ & 3 \\ \hline
$10 \vdash 20$ & 5 \\ \hline
$20 \vdash 30$ & 8 \\ \hline
$30 \vdash 40$ & 4 \\ \hline
$40 \vdash 50$ & 2 \\ \hline
\end{tabular}
\end{center}

Com base nessa tabela, classifique os dados quanto sua assimetria, usando o 2.\textordmasculine\ coeficiente de assimetria de Pearson, \(e_2\). Para isso, considere a seguinte classificação:
\begin{center}
    \renewcommand{\arraystretch}{1.2}
    \begin{tabular}{|c|c|}
        \hline
        \textbf{Valor de } $e_2$ & \textbf{Classificação da Assimetria} \\ \hline
        $|e_2| = 0$ & Simétrica \\ \hline
        $0 < |e_2| \le 0{,}15$ & Assimetria fraca \\ \hline
        $0{,}15 < |e_2| \le 1$ & Assimetria moderada \\ \hline
        $|e_2| > 1$ & Assimetria forte \\ \hline
    \end{tabular}
\end{center}
\end{exercicioBanco}
